\chapter{Vectors in $\R^n$}

\section{Vectors and their operations}

\begin{definition}[$\R^n$]
$\R^n$ is the set of all ordered $n$-tuples of real numbers, written as column vectors
\[
  \vec{x} = \begin{bmatrix} x_1 \\ x_2 \\ \vdots \\ x_n \end{bmatrix}, \qquad x_1,\dots,x_n \in \R .
\]
The number $x_i$ is the $i$-th \emph{component} (or entry) of $\vec x$. Two vectors are equal
when all their components agree. The \emph{zero vector} is $\vec 0 = (0,\dots,0)^T$.
\end{definition}

To save space we sometimes write $\vec x = (x_1,\dots,x_n)$ or $\vec x = [x_1 \ \cdots \ x_n]^T$,
but in this course vectors are \emph{column} vectors.

\begin{definition}[Addition and scalar multiplication]
For $\vec x,\vec y\in\R^n$ and $c\in\R$,
\[
  \vec x + \vec y = \begin{bmatrix} x_1+y_1 \\ \vdots \\ x_n+y_n\end{bmatrix}, \qquad
  c\vec x = \begin{bmatrix} cx_1 \\ \vdots \\ cx_n \end{bmatrix}.
\]
We write $-\vec x = (-1)\vec x$ and $\vec x - \vec y = \vec x + (-\vec y)$.
\end{definition}

\begin{theorem}[Properties of $\R^n$]\label{thm:rn-props}
For all $\vec x,\vec y,\vec w\in\R^n$ and $c,d\in\R$:
\begin{enumerate}[label=V\arabic*.]
  \item $\vec x+\vec y\in\R^n$ (closed under addition);
  \item $(\vec x+\vec y)+\vec w = \vec x+(\vec y+\vec w)$;
  \item $\vec x+\vec y = \vec y+\vec x$;
  \item $\vec x+\vec 0 = \vec x$;
  \item $\vec x+(-\vec x) = \vec 0$;
  \item $c\vec x\in\R^n$ (closed under scalar multiplication);
  \item $c(d\vec x) = (cd)\vec x$;
  \item $(c+d)\vec x = c\vec x+d\vec x$;
  \item $c(\vec x+\vec y) = c\vec x+c\vec y$;
  \item $1\vec x = \vec x$.
\end{enumerate}
\end{theorem}
\begin{proof}
Each follows componentwise from the corresponding property of real numbers. For example
for V8, the $i$-th component of $(c+d)\vec x$ is $(c+d)x_i = cx_i + dx_i$, which is the
$i$-th component of $c\vec x+d\vec x$.
\end{proof}

These ten properties will later become the \emph{axioms} of an abstract vector space
(Chapter 11).

\section{Linear combinations and span}

\begin{definition}[Linear combination]
For vectors $\vec v_1,\dots,\vec v_k\in\R^n$ and scalars $c_1,\dots,c_k\in\R$, the vector
$c_1\vec v_1+\dots+c_k\vec v_k$ is a \emph{linear combination} of $\vec v_1,\dots,\vec v_k$.
\end{definition}

\begin{definition}[Span]
The \emph{span} of $B=\{\vec v_1,\dots,\vec v_k\}$ is the set of all linear combinations:
\[
  \Span B = \{c_1\vec v_1+\dots+c_k\vec v_k : c_1,\dots,c_k\in\R\}.
\]
If $S=\Span B$ we say $B$ \emph{spans} $S$ and that $B$ is a \emph{spanning set} for $S$.
By convention $\Span\emptyset = \{\vec 0\}$.
\end{definition}

\begin{example}
Is $\vec x = (1,2,3)$ in $\Span\{(1,0,1),(0,1,1)\}$? We need
$c_1(1,0,1)+c_2(0,1,1) = (c_1, c_2, c_1+c_2) = (1,2,3)$. The first two components force
$c_1=1$, $c_2=2$, and then $c_1+c_2=3$ agrees with the third component. So yes:
$\vec x = 1(1,0,1)+2(0,1,1)$. On the other hand $(1,2,4)$ is \emph{not} in the span, since
$c_1+c_2=3\ne 4$.
\end{example}

\begin{theorem}[Removing redundant vectors]\label{thm:span-reduce}
Let $\vec v_1,\dots,\vec v_k\in\R^n$. If $\vec v_k$ is a linear combination of
$\vec v_1,\dots,\vec v_{k-1}$, then
$\Span\{\vec v_1,\dots,\vec v_k\} = \Span\{\vec v_1,\dots,\vec v_{k-1}\}$.
\end{theorem}
\begin{proof}
The inclusion $\supseteq$ is clear (take coefficient $0$ on $\vec v_k$). Conversely, write
$\vec v_k = a_1\vec v_1+\dots+a_{k-1}\vec v_{k-1}$. Then any
$c_1\vec v_1+\dots+c_k\vec v_k = (c_1+c_ka_1)\vec v_1+\dots+(c_{k-1}+c_ka_{k-1})\vec v_{k-1}$
lies in the smaller span.
\end{proof}

\section{Linear independence}

\begin{definition}[Linear (in)dependence]
A set $\{\vec v_1,\dots,\vec v_k\}\subseteq\R^n$ is \emph{linearly dependent} if there are
scalars $c_1,\dots,c_k$, \emph{not all zero}, with $c_1\vec v_1+\dots+c_k\vec v_k=\vec 0$.
Otherwise it is \emph{linearly independent}: the only solution of
$c_1\vec v_1+\dots+c_k\vec v_k=\vec 0$ is the \emph{trivial solution} $c_1=\dots=c_k=0$.
\end{definition}

\begin{theorem}
A set $\{\vec v_1,\dots,\vec v_k\}$ is linearly dependent if and only if some $\vec v_i$ is a
linear combination of the others. In particular, any set containing $\vec 0$ is linearly
dependent.
\end{theorem}
\begin{proof}
If $c_1\vec v_1+\dots+c_k\vec v_k=\vec 0$ with $c_i\ne0$, then
$\vec v_i = -\frac{1}{c_i}\sum_{j\ne i}c_j\vec v_j$. Conversely, if
$\vec v_i=\sum_{j\ne i}a_j\vec v_j$, then $\sum_{j\ne i}a_j\vec v_j - 1\vec v_i=\vec 0$ is a
nontrivial relation (coefficient of $\vec v_i$ is $-1\neq 0$). If $\vec v_1=\vec 0$, then
$1\vec v_1 + 0\vec v_2+\dots=\vec 0$.
\end{proof}

\begin{example}
Is $\{(1,2,1),(0,1,1),(1,3,2)\}$ linearly independent? Consider
$c_1(1,2,1)+c_2(0,1,1)+c_3(1,3,2)=\vec 0$, i.e.
\[
  c_1+c_3=0,\qquad 2c_1+c_2+3c_3=0,\qquad c_1+c_2+2c_3=0.
\]
From the first, $c_1=-c_3$; substituting into the third gives $c_2=-c_3$; the second
becomes $-2c_3-c_3+3c_3=0$, which always holds. Taking $c_3=1$ gives the nontrivial
relation $-(1,2,1)-(0,1,1)+(1,3,2)=\vec 0$, so the set is linearly \emph{dependent}. Indeed
$(1,3,2)=(1,2,1)+(0,1,1)$.
\end{example}

In Chapter 3 we will solve such systems systematically by row reduction.

\section{Bases of $\R^n$ and subspaces}

\begin{definition}[Basis]
A set $B$ of vectors in a subspace $S$ (or in $\R^n$) is a \emph{basis} for $S$ if $B$ is
linearly independent and $\Span B = S$.
\end{definition}

\begin{example}[Standard basis]
The \emph{standard basis} of $\R^n$ is $\{\vec e_1,\dots,\vec e_n\}$, where $\vec e_i$ has a $1$ in
position $i$ and $0$ elsewhere. Every $\vec x$ is $x_1\vec e_1+\dots+x_n\vec e_n$, and
$c_1\vec e_1+\dots+c_n\vec e_n=(c_1,\dots,c_n)=\vec 0$ forces all $c_i=0$.
\end{example}

\begin{theorem}[Unique representation]
If $B=\{\vec v_1,\dots,\vec v_k\}$ is a basis for $S$, then every $\vec x\in S$ can be written
as a linear combination of the vectors of $B$ in \emph{exactly one} way.
\end{theorem}
\begin{proof}
Existence holds because $B$ spans $S$. If $\vec x=\sum c_i\vec v_i=\sum d_i\vec v_i$, subtracting
gives $\sum(c_i-d_i)\vec v_i=\vec 0$, so $c_i=d_i$ for all $i$ by independence.
\end{proof}

We will see (Chapter 3) that a set of $k$ vectors in $\R^n$ can only be linearly independent
if $k\le n$ and can only span $\R^n$ if $k\ge n$; hence every basis of $\R^n$ has exactly $n$
vectors, and \emph{a set of $n$ vectors in $\R^n$ is a basis iff it is linearly independent iff
it spans $\R^n$.}

\begin{definition}[Subspace]
A subset $S\subseteq\R^n$ is a \emph{subspace} of $\R^n$ if $S$ together with the
operations of $\R^n$ satisfies all of V1--V10.
\end{definition}

Most of the properties are inherited automatically from $\R^n$, which gives a short test.

\begin{theorem}[Subspace Test]
A subset $S\subseteq\R^n$ is a subspace if and only if
\begin{enumerate}
  \item $S$ is non-empty (equivalently, $\vec 0\in S$),
  \item $\vec x+\vec y\in S$ for all $\vec x,\vec y\in S$, and
  \item $c\vec x\in S$ for all $\vec x\in S$, $c\in\R$.
\end{enumerate}
\end{theorem}
\begin{proof}[Proof sketch]
V2, V3, V7--V10 hold for vectors in $S$ because they hold in $\R^n$. V1 and V6 are
conditions 2 and 3. For V4 and V5: pick $\vec x\in S$; then $0\vec x=\vec 0\in S$ and
$(-1)\vec x=-\vec x\in S$ by condition 3.
\end{proof}

\begin{example}
$S=\{\vec x\in\R^3 : x_1+x_2-x_3=0\}$ is a subspace. (1) $0+0-0=0$, so $\vec 0\in S$.
(2) If $\vec x,\vec y\in S$ then $(x_1+y_1)+(x_2+y_2)-(x_3+y_3) = (x_1+x_2-x_3)+(y_1+y_2-y_3)=0$.
(3) $cx_1+cx_2-cx_3=c(x_1+x_2-x_3)=0$. Solving $x_3=x_1+x_2$ shows
$S=\Span\{(1,0,1),(0,1,1)\}$, and these two vectors are clearly independent, so they form a
basis of $S$.
\end{example}

\begin{example}
$T=\{\vec x\in\R^2 : x_1=1\}$ is \emph{not} a subspace since $\vec 0\notin T$. And
$U=\{\vec x\in\R^2 : x_1x_2=0\}$ (the two axes) is not a subspace: $(1,0),(0,1)\in U$ but
$(1,1)\notin U$.
\end{example}

\begin{theorem}
For any $\vec v_1,\dots,\vec v_k\in\R^n$, $\Span\{\vec v_1,\dots,\vec v_k\}$ is a subspace of $\R^n$.
\end{theorem}
\begin{proof}
$\vec 0 = 0\vec v_1+\dots+0\vec v_k$ is in the span; sums and scalar multiples of linear
combinations are again linear combinations.
\end{proof}

\section{Lines and planes}

\begin{definition}
Let $\vec p,\vec d\in\R^n$ with $\vec d\neq\vec 0$. The set with \emph{vector equation}
$\vec x = \vec p + t\vec d$, $t\in\R$, is the \emph{line} through $\vec p$ with direction
$\vec d$. If $\{\vec u,\vec v\}$ is linearly independent, $\vec x=\vec p+s\vec u+t\vec v$,
$s,t\in\R$, is the \emph{plane} through $\vec p$ spanned by $\vec u,\vec v$.
\end{definition}

Geometrically, the span of one non-zero vector is a line through the origin, and the span
of two independent vectors is a plane through the origin. A line or plane that does not
pass through the origin is a \emph{translate} of a subspace, not a subspace itself.

\begin{example}
The line through $P(1,2,-1)$ and $Q(3,0,4)$ has direction $\vec{PQ}=(2,-2,5)$, so a vector
equation is $\vec x = (1,2,-1) + t(2,-2,5)$, $t\in\R$. Its \emph{parametric equations} are
$x_1=1+2t$, $x_2=2-2t$, $x_3=-1+5t$.
\end{example}

\section{Dot product, length and angles}

\begin{definition}[Dot product, length]
For $\vec x,\vec y\in\R^n$, the \emph{dot product} is
$\vec x\cdot\vec y = x_1y_1+\dots+x_ny_n$, and the \emph{length} (norm) of $\vec x$ is
$\norm{\vec x}=\sqrt{\vec x\cdot\vec x}$. A vector of length $1$ is a \emph{unit vector};
for $\vec x\neq\vec 0$, $\hat x = \frac{1}{\norm{\vec x}}\vec x$ is the unit vector in the
direction of $\vec x$.
\end{definition}

\begin{theorem}[Properties of the dot product]
For $\vec x,\vec y,\vec z\in\R^n$, $c\in\R$:
(1) $\vec x\cdot\vec x\ge0$ with equality iff $\vec x=\vec 0$;
(2) $\vec x\cdot\vec y=\vec y\cdot\vec x$;
(3) $\vec x\cdot(\vec y+\vec z)=\vec x\cdot\vec y+\vec x\cdot\vec z$;
(4) $(c\vec x)\cdot\vec y = c(\vec x\cdot\vec y)$;
(5) $\norm{c\vec x}=\abs{c}\norm{\vec x}$;
(6) $\abs{\vec x\cdot\vec y}\le\norm{\vec x}\norm{\vec y}$ (Cauchy--Schwarz);
(7) $\norm{\vec x+\vec y}\le\norm{\vec x}+\norm{\vec y}$ (triangle inequality).
\end{theorem}
\begin{proof}[Proof of (6) and (7)]
If $\vec y=\vec 0$ both sides of (6) are $0$. Otherwise, for every $t\in\R$,
\[
0\le \norm{\vec x+t\vec y}^2 = \norm{\vec y}^2t^2 + 2(\vec x\cdot\vec y)t + \norm{\vec x}^2 .
\]
A quadratic in $t$ that is never negative has discriminant $\le0$:
$4(\vec x\cdot\vec y)^2-4\norm{\vec x}^2\norm{\vec y}^2\le0$, which is (6). For (7),
$\norm{\vec x+\vec y}^2=\norm{\vec x}^2+2\vec x\cdot\vec y+\norm{\vec y}^2\le
\norm{\vec x}^2+2\norm{\vec x}\norm{\vec y}+\norm{\vec y}^2=(\norm{\vec x}+\norm{\vec y})^2$.
\end{proof}

By Cauchy--Schwarz, $-1\le\frac{\vec x\cdot\vec y}{\norm{\vec x}\norm{\vec y}}\le1$ for non-zero
vectors, so we can \emph{define} the angle $\theta\in[0,\pi]$ between them by
\[
  \cos\theta = \frac{\vec x\cdot\vec y}{\norm{\vec x}\,\norm{\vec y}}.
\]
In $\R^2$ and $\R^3$ this agrees with the usual angle (law of cosines).

\begin{example}
Let $\vec u=(1,2,2)$, $\vec v=(3,0,4)$. Then $\vec u\cdot\vec v = 3+0+8=11$,
$\norm{\vec u}=\sqrt{1+4+4}=3$, $\norm{\vec v}=\sqrt{9+0+16}=5$, so
$\cos\theta = 11/15$ and $\theta = \arccos(11/15)\approx 0.75$ rad.
\end{example}

\begin{definition}[Orthogonal]
$\vec x,\vec y\in\R^n$ are \emph{orthogonal} if $\vec x\cdot\vec y=0$. (So $\vec 0$ is orthogonal
to everything.) A set of vectors is \emph{orthogonal} if every pair of distinct vectors in it is
orthogonal.
\end{definition}

\begin{definition}[Scalar equation of a plane]
In $\R^3$, a non-zero vector $\vec n$ is a \emph{normal vector} of a plane if it is orthogonal to
every direction vector in the plane. If the plane passes through $P$, then $X$ is on the plane
iff $\vec n\cdot(\vec x-\vec p)=0$. With $\vec n=(a,b,c)$ this is the scalar equation
$ax_1+bx_2+cx_3 = \vec n\cdot\vec p$. More generally $\vec n\cdot\vec x = d$ in $\R^n$ defines a
\emph{hyperplane}.
\end{definition}

\section{The cross product in $\R^3$}

\begin{definition}[Cross product]
For $\vec u=(u_1,u_2,u_3)$, $\vec v=(v_1,v_2,v_3)$ in $\R^3$,
\[
  \vec u\times\vec v = \begin{bmatrix} u_2v_3-u_3v_2 \\ u_3v_1-u_1v_3 \\ u_1v_2-u_2v_1\end{bmatrix}.
\]
\end{definition}

\begin{theorem}
For $\vec u,\vec v,\vec w\in\R^3$ and $c\in\R$:
\begin{enumerate}
  \item $\vec u\times\vec v$ is orthogonal to both $\vec u$ and $\vec v$;
  \item $\vec v\times\vec u = -(\vec u\times\vec v)$, so $\vec u\times\vec u=\vec 0$;
  \item $\vec u\times\vec v=\vec 0$ iff one of $\vec u,\vec v$ is a scalar multiple of the other;
  \item $\vec u\times(\vec v+\vec w)=\vec u\times\vec v+\vec u\times\vec w$ and
        $(c\vec u)\times\vec v = c(\vec u\times\vec v)$;
  \item $\norm{\vec u\times\vec v}=\norm{\vec u}\norm{\vec v}\sin\theta$, which is the area of the
        parallelogram spanned by $\vec u$ and $\vec v$.
\end{enumerate}
\end{theorem}
\begin{proof}[Proof of (1)]
$\vec u\cdot(\vec u\times\vec v)= u_1(u_2v_3-u_3v_2)+u_2(u_3v_1-u_1v_3)+u_3(u_1v_2-u_2v_1)$;
the six terms cancel in pairs. Similarly for $\vec v$.
\end{proof}

\begin{important}[Pitfall]
The cross product is \emph{not} associative: $\vec e_1\times(\vec e_1\times\vec e_2) =
\vec e_1\times\vec e_3 = -\vec e_2$, but $(\vec e_1\times\vec e_1)\times\vec e_2=\vec 0$. It is
only defined in $\R^3$.
\end{important}

\begin{example}
$(1,2,3)\times(4,5,6) = (2\cdot6-3\cdot5,\ 3\cdot4-1\cdot6,\ 1\cdot5-2\cdot4) = (-3,6,-3)$.
Check: $(1,2,3)\cdot(-3,6,-3)=-3+12-9=0$ and $(4,5,6)\cdot(-3,6,-3)=-12+30-18=0$.
\end{example}

\begin{example}[Plane through three points]
Find the scalar equation of the plane through $P(1,0,0)$, $Q(0,2,0)$, $R(0,0,3)$. Direction
vectors are $\vec{PQ}=(-1,2,0)$ and $\vec{PR}=(-1,0,3)$, so a normal is
\[
\vec n=\vec{PQ}\times\vec{PR}=\big(2\cdot3-0\cdot0,\ 0\cdot(-1)-(-1)\cdot3,\ (-1)\cdot0-2\cdot(-1)\big)=(6,3,2).
\]
The plane is $6x_1+3x_2+2x_3 = \vec n\cdot\vec p = 6$. Check: $Q$ gives $6$, $R$ gives $6$.
The triangle $PQR$ has area $\frac12\norm{\vec n}=\frac12\sqrt{36+9+4}=\frac72$.
\end{example}

\section*{Summary}
\begin{itemize}
  \item $\R^n$ with componentwise $+$ and scalar multiplication satisfies ten properties V1--V10.
  \item $\Span$ = all linear combinations; it is always a subspace. A redundant vector can be
        removed without changing the span.
  \item Linearly independent = only the trivial relation; dependent iff some vector is a
        combination of the others. A basis is an independent spanning set and gives unique
        coordinates.
  \item Subspace test: contains $\vec 0$, closed under $+$ and scalar multiplication.
  \item $\vec x\cdot\vec y$, $\norm{\vec x}$, $\cos\theta=\frac{\vec x\cdot\vec y}{\norm{\vec x}\norm{\vec y}}$;
        Cauchy--Schwarz and triangle inequality; orthogonal means dot product $0$.
  \item $\vec u\times\vec v$ is orthogonal to both, has length equal to the parallelogram's area,
        and gives normal vectors of planes in $\R^3$.
\end{itemize}
